
\documentclass[11pt]{article}

\usepackage[margin=1in]{geometry}
\usepackage{amsmath,amssymb}
\usepackage{bm}
\usepackage{enumitem}
\usepackage{tikz}
\usetikzlibrary{arrows.meta,positioning}

\title{Surface-Related Multiples as a Feedback System}
\author{}
\date{}

\begin{document}

\maketitle

\section{Introduction}

Surface-related multiples can be understood naturally using the language of
\emph{feedback systems}. The idea is simple: a primary reflection reaches the
surface, is reflected back downward, interacts with the subsurface again, and
therefore generates additional reflected arrivals. Repetition of this process
creates higher-order multiples.

We use the following geophysical notation:

\begin{itemize}[leftmargin=2em]
    \item $w$: source wavelet,
    \item $P$: primary reflection response, or primary reflection operator,
    \item $R$: reflection coefficient at the surface,
    \item $M$: multiple contribution to the data.
\end{itemize}

Throughout this note, $*$ denotes convolution.

\section{Primaries Without Feedback}

If there were no reflection at the free surface, the source wavelet $w$ would
propagate into the earth, interact with the subsurface reflectors, and generate
the primary reflection response

\begin{equation}
d_P = P * w.
\end{equation}

Thus, $P$ represents the action of the subsurface on the incident wavelet.

\section{The Surface as a Feedback Mechanism}

When the primary wavefield reaches the surface, it is not necessarily removed
from the system. Instead, it can be reflected back downward with reflection
coefficient $R$.

The reflected wave then propagates through the earth again and interacts once
more with the primary reflection operator $P$.

This produces the first-order multiple:

\begin{equation}
M_1 = P * R * P * w.
\end{equation}

The physical sequence is therefore

\begin{equation}
w
\;\longrightarrow\;
P
\;\longrightarrow\;
R
\;\longrightarrow\;
P.
\end{equation}

The first application of $P$ creates the primary reflection. The factor $R$
represents reflection at the surface, and the second application of $P$
represents another interaction with the subsurface.

\section{Higher-Order Multiples}

After the first-order multiple reaches the surface, it can again be reflected
downward. This produces a second-order multiple,

\begin{equation}
M_2 = P * R * P * R * P * w.
\end{equation}

Similarly, the third-order multiple is

\begin{equation}
M_3 = P * R * P * R * P * R * P * w.
\end{equation}

Each additional trip through the feedback loop introduces another factor
$R * P$.

The total recorded wavefield can therefore be written as

\begin{equation}
d =
P*w
+
P*R*P*w
+
P*R*P*R*P*w
+\cdots .
\end{equation}

The first term is the primary contribution, while the remaining terms are
surface-related multiples.

Defining the total multiple contribution as $M$,

\begin{equation}
M =
P*R*P*w
+
P*R*P*R*P*w
+\cdots .
\end{equation}

Hence,

\begin{equation}
d = P*w + M.
\end{equation}

\section{Operator Form}

If convolution symbols are omitted for compactness, the data can be written as

\begin{equation}
d = Pw + PRPw + PRPRPw + \cdots .
\end{equation}

Factoring out the first primary operator gives

\begin{equation}
d =
P
\left[
I + RP + (RP)^2 + (RP)^3 + \cdots
\right]w.
\end{equation}

The expression inside brackets is a geometric operator series. Formally, when
the series converges,

\begin{equation}
I + RP + (RP)^2 + \cdots = (I-RP)^{-1}.
\end{equation}

Therefore,

\begin{equation}
\boxed{
d = P(I-RP)^{-1}w
}
\end{equation}

which has the characteristic form of a system with feedback.

\section{Feedback Interpretation}

The key physical interpretation is that the primary wavefield is
\emph{fed back} into the earth after reflecting at the surface.

A simplified feedback representation is shown in Figure~\ref{fig:feedback}.

\begin{figure}[h]
\centering
\begin{tikzpicture}[
    node distance=1.8cm,
    block/.style={draw, rectangle, minimum width=1.5cm, minimum height=0.9cm},
    sum/.style={draw, circle, inner sep=1.5pt},
    >=Latex
]

\node (w) {$w$};
\node[block, right=of w] (P) {$P$};
\node[sum, right=of P] (sum) {$+$};
\node[right=of sum] (d) {$d$};

\draw[->] (w) -- (P);
\draw[->] (P) -- (sum);
\draw[->] (sum) -- (d);

\node[block, below=1.4cm of sum] (R) {$R$};
\node[block, left=of R] (Pf) {$P$};

\draw[->] (d) |- (R);
\draw[->] (R) -- (Pf);
\draw[->] (Pf) -| (sum);

\end{tikzpicture}
\caption{Conceptual representation of surface-related multiples as a feedback
system. The wavefield reaching the surface is reflected by $R$ and re-enters
the earth, where it interacts again with the primary reflection operator $P$.}
\label{fig:feedback}
\end{figure}

Each traversal of the feedback path generates one additional order of multiple.

\section{Role of the Surface Reflection Coefficient}

For pressure waves at an ideal free surface, the reflection coefficient is
often approximately

\begin{equation}
R \approx -1.
\end{equation}

The negative sign represents the polarity reversal associated with reflection
of pressure at the free surface.

Thus, the surface does more than merely return energy into the subsurface: it
can also change the phase or polarity of that returned wavefield.

\section{Main Interpretation}

The feedback view provides a compact way to understand multiple generation:

\begin{quote}
A primary reflection reaches the surface. The surface reflects it back into the
earth. The returned wavefield interacts again with the subsurface reflectors,
creating a multiple. Repetition of this feedback process creates higher-order
multiples.
\end{quote}

In this interpretation,

\begin{equation}
\boxed{
\text{primaries become the input to the feedback loop that generates multiples.}
}
\end{equation}

This viewpoint is useful because it connects geophysical multiple prediction
with the general theory of systems with feedback. If the primary reflection
operator $P$ and the surface reflection operator $R$ can be estimated, the
multiple wavefield can in principle be predicted by repeatedly applying the
feedback process.

\end{document}
